\uuid{Szzc}
\titre{Intégration par parties}
\auteur{hanani ; mzari}
\organisation{exo7}
\licence{CC BY-NC-SA 4.0}
\datecreate{2020-01-20}
\niveau{L1}
\module{Analyse}
\chapitre{Calcul d'intégrales}
\sousChapitre{Autre}
\difficulte{1}
\theme{Calculs d'intégrales, intégration par parties}
\exo7qcmid{lille-2.integrales.005}

\contenu{
\begin{question}
\qclassification{127.99}
\qtags{motcle=intégration par parties}
Parmi les affirmations suivantes, cocher celles qui sont vraies :
\begin{answers}
\good{$\displaystyle \int _1^{\mathrm{e}}\ln t\mathrm{d}t=\Big[t\ln t\Big]_1^{\mathrm{e}}-\int _1^{\mathrm{e}}\mathrm{d}t$.}
\good{$\displaystyle \int _1^{\mathrm{e}}\ln t\mathrm{d}t=1$.}
\bad{$\displaystyle \int _1^{2}t\ln t\mathrm{d}t=\Big[t^2\ln t\Big]_1^{\mathrm{e}}-\int _1^{\mathrm{e}}t\mathrm{d}t$.}
\bad{$\displaystyle \int _1^{2}t\ln t\mathrm{d}t=\frac{\mathrm{e}^2+1}{2}$.}
\end{answers}
\vskip2mm
\begin{explanations}
Une intégration par parties avec $u=\ln t$ et $v=t\Rightarrow v'=1$ donne :
$$\displaystyle \int _1^{\mathrm{e}}\ln t\mathrm{d}t=\Big[t\ln t\Big]_1^{\mathrm{e}}-\int _1^{\mathrm{e}}\mathrm{d}t=1.$$
Une intégration par parties, avec $u=\ln t$ et $v=t^2/2\Rightarrow v'=t$, donne :
$$\displaystyle \int _1^{2}t\ln t\mathrm{d}t=\left[\frac{t^2}{2}\ln t\right]_1^{\mathrm{e}}-\int _1^{\mathrm{e}}\frac{t}{2}\mathrm{d}t=\frac{\mathrm{e}^2+1}{4}.$$
\end{explanations}
\end{question}
}
